\documentclass[12pt]{article}
\usepackage[utf8]{inputenc}
\usepackage[T1]{fontenc}
\usepackage[margin=1in]{geometry}
\usepackage{setspace}
\usepackage{xcolor}
\usepackage[colorlinks=true, urlcolor=blue!60!black, linkcolor=blue!60!black]{hyperref}
\setlength{\parindent}{2em}
\setlength{\parskip}{1.25em}
\renewcommand{\baselinestretch}{1.0}
\begin{document}
\begin{center}
{\LARGE \textit{Moore v. Harper}, 600 U.S.\ 1 (2023)\par}
\vspace{0.8em}
{\large Case Brief --- Four-Page Edition\par}
\end{center}
\vspace{1em}

\section{Memory Jogger}

The Elections Clause does not insulate state legislatures from ordinary state-court judicial review under state constitutions when they regulate federal elections, though federal courts retain a duty to ensure state courts do not transgress the ordinary bounds of judicial review and arrogate legislative power (\href{https://www.law.cornell.edu/supremecourt/text/21-1271}{Cornell LII}).

\section{Detailed Case Facts}

After the 2020 census gave North Carolina an additional House seat, the Republican General Assembly enacted new state legislative and congressional maps along party lines in 2021. Voters including Rebecca Harper, the League of Conservation Voters, and Common Cause sued House Speaker Timothy Moore and others in state court, alleging extreme partisan gerrymanders violating the state constitution's Free Elections Clause (``[a]ll elections shall be free''), Equal Protection, Free Speech, and Free Assembly Clauses. Defending the congressional map, the legislators argued that gerrymandering claims were nonjusticiable and, under the ``independent state legislature theory'' (ISLT), that the Federal Elections Clause gives state legislatures exclusive, unchecked authority over congressional maps.

\section{Procedural History}

A three-judge trial court found the map ``a partisan outlier intentionally and carefully designed to maximize Republican advantage'' but dismissed the claims as nonjusticiable. In \textit{Harper I} (2022), the North Carolina Supreme Court reversed, distinguished \textit{Rucho v. Common Cause}, rejected the Elections Clause defense, and enjoined the maps; the legislature's remedial map included a ``trigger'' reviving the 2021 map if \textit{Harper I} were reversed. The U.S. Supreme Court granted certiorari. After the 2022 elections gave Republicans a 5--2 state-court majority, \textit{Harper III} (2023) overruled \textit{Harper I} and held gerrymandering nonjusticiable, without reinstating the 2021 maps. The U.S. Supreme Court decided the case on June 27, 2023.

\section{Judicial Votes}

Affirmed, \textbf{6--3}.
\begin{itemize}
\item \textbf{Majority:} Chief Justice Roberts, joined by Justices Sotomayor, Kagan, Kavanaugh, Barrett, and Jackson.
\item \textbf{Concurring:} Justice Kavanaugh.
\item \textbf{Dissenting:} Justice Thomas, joined by Justice Gorsuch, and by Justice Alito as to Part~I (mootness).
\end{itemize}

\section{Holding}

The case was not moot, and \textit{Harper I} was final under 28 U.S.C.\ \S~1257(a), because \textit{Harper III} did not vacate its injunction and the trigger would revive the 2021 maps on reversal. On the merits, the Elections Clause does not vest state legislatures with authority free from state constitutional constraints or state judicial review, with no procedure-versus-substance carve-out. But state courts ``do not have free rein''; federal courts must ensure they do not exceed the ordinary bounds of judicial review. The Court adopted no test for that limit, a discussion commentators call ``ISLT-lite.''

\section{Analysis of Opinions}

\textbf{Majority (Roberts).} Founding-era state decisions such as North Carolina's \textit{Bayard v. Singleton} show judicial review was a ``fundamental principle'' the Elections Clause did not silently displace. \textit{Hildebrant}, \textit{Smiley v. Holm}, and \textit{Arizona State Legislature} treat election rules as lawmaking subject to state constitutions; legislatures are ``creatures of the State Constitutions.'' The procedure-substance line is ``hazy,'' and early state constitutions routinely regulated federal elections. Drawing on Takings and Contracts Clause cases and the \textit{Bush v. Gore} opinions, federal courts may police state courts' outer bounds, but petitioners disclaimed any argument that \textit{Harper I} misread state law.

\textbf{Kavanaugh (concurring).} He endorsed Chief Justice Rehnquist's \textit{Bush v. Gore} standard---whether a state court ``impermissibly distorted'' state law ``beyond what a fair reading required''---for both statutes and constitutions, measured against state law as it existed before the decision.

\textbf{Thomas (dissenting).} The case was ``indisputably moot'' and the opinion ``plainly advisory,'' since \textit{Harper III} gave petitioners complete relief. On the merits (with Gorsuch), state constitutions may define who makes laws and how, but not impose substantive limits on what federal election rules the legislature enacts. Federal review of open-textured state constitutions would federalize state law without intelligible standards.

\section{Examples: Future Applications}

\begin{enumerate}
\item \textbf{Same-Side: Express Compactness Rules.} A state court strikes a map violating the state constitution's compactness and county-boundary text. That is ordinary judicial review.
\item \textbf{Same-Side: Settled Precedent on Drop Boxes.} Relying on decades of precedent under a ``free and equal elections'' clause, a state court strikes a drop-box restriction. It stays within ordinary bounds.
\item \textbf{Opposite-Side: Rewriting a Deadline.} Days before an election, a court extends an unambiguous ballot deadline by two weeks on ``policy fairness'' alone. A federal court could intervene.
\item \textbf{Opposite-Side: Map Without a Violation.} A state court imposes an expert's map without finding the enacted map unlawful, intruding on ``the Legislature'' under Article I, \S~4.
\item \textbf{Fence-Sitter: Emergency Extension.} During a hurricane, a court reads a broad ``free elections'' guarantee for the first time to extend a mail-ballot deadline by three days. Defenders call it common-law judging; challengers invoke Kavanaugh's pre-existing-law baseline.
\end{enumerate}

\section{Critique}

\textbf{Originalist views.} Akhil Amar, Vikram Amar, and Steven Calabresi celebrated the rejection of ISLT, arguing that public meaning, the Articles of Confederation, and uniform early state constitutions refute it and that the ``Pinckney Plan'' fails historical scrutiny. William Baude and Michael McConnell questioned whether state courts may draw replacement maps, and formalists echoed Thomas that the Court issued an advisory opinion.

\textbf{Progressive views.} Richard Hasen, Richard Pildes, Carolyn Shapiro, Leah Litman, and Katherine Shaw warned that Part V planted a ``time bomb'': the undefined ``ordinary bounds'' standard and Kavanaugh's anti-novelty gloss invite federal judges to impose their methodology on state courts. The \textit{Harvard Law Review} noted that the Court curbed legislatures only by making federal judges supervisors of state election law, reviving \textit{Bush v. Gore} in a majority opinion.

\section{Key Quotations}

\begin{quote}
``The Elections Clause does not insulate state legislatures from the ordinary exercise of state judicial review.''
\end{quote}
\noindent --- Chief Justice Roberts, Opinion of the Court.

\begin{quote}
``We do not adopt these or any other test by which we can measure state court interpretations of state law in cases implicating the Elections Clause. The questions presented in this area are complex and context specific. We hold only that state courts may not transgress the ordinary bounds of judicial review such that they arrogate to themselves the power vested in state legislatures to regulate federal elections.''
\end{quote}
\noindent --- Chief Justice Roberts, Opinion of the Court.

\begin{quote}
``The legislature acts both as a lawmaking body created and bound by its state constitution, and as the entity assigned particular authority by the Federal Constitution. Both constitutions restrain the legislature's exercise of power.''
\end{quote}
\noindent --- Chief Justice Roberts, Opinion of the Court.

\end{document}
